\exercisesfor{3}

\begin{enumerate}
\def\labelenumi{\arabic{enumi}.}
\item
  Let g be a Lie algebra, \(\rho\) a representation of g on a K-module M and \(\sigma\) the associated representation of g on the tensor algebra of M. Show that the submodule of symmetric tensors and the submodule of skew-symmetric tensors are stable under \(\sigma\) .
\item
  Let g be a Lie algebra, \(\rho\) a representation of g on a K-module M and \(\sigma\) the associated representation of g on Q = L(M, M; M). For \(f \in Q\) to be invariant under \(\sigma\) , it is necessary and sufficient that the \(\rho(x)\) be derivations of M with the multiplication defined by f.
\item
  Let V be a finite-dimensional vector space over a perfect field K.
\end{enumerate}

\begin{enumerate}
\def\labelenumi{(\alph{enumi})}
\tightlist
\item
  The identity representation of \(\mathfrak{gl}(\mathrm{V})\) defines canonically representations of \(\mathfrak{gl}(\mathrm{V})\) on \(\bigotimes_{q}^{p}\mathrm{V}\) and \(\bigotimes_{q}^{q}\mathrm{V}^{*}\) and hence a representation \(u\mapsto u_q^p\) of \(\mathfrak{gl}(\mathrm{V})\) on
\end{enumerate}

\(V_{q}^{p} = \left( \bigotimes_{i=1}^{p} V \right) \otimes \left( \bigotimes_{i=1}^{q} V^{*} \right)\) . Show that \(u_{q}^{p}\) is semi-simple if \(u\) is. (Reduce it, by extending the base field, to the case where \(K\) is algebraically closed.) Show that \(u_{q}^{p}\) is nilpotent if \(u\) is. Deduce that, if \(s\) and \(n\) are the semi-simple and nilpotent components of \(u, u_{q}^{p}\) and \(n_{q}^{p}\) are the semi-simple and nilpotent components of \(u_{q}^{p}\) .

\begin{enumerate}
\def\labelenumi{(\alph{enumi})}
\setcounter{enumi}{1}
\item
  Deduce from (a) that, if an element of \(V_q^p\) is annihilated by \(u_q^p\) , it is annihilated by \(s_q^p\) and \(n_q^p\) .
\item
  Deduce from (b) and Exercise 2 that, if V has a not necessarily associative algebra structure and u is a derivation of V, then s and n are derivations of V.
\end{enumerate}

\begin{enumerate}
\def\labelenumi{\arabic{enumi}.}
\setcounter{enumi}{3}
\tightlist
\item
  Suppose that K is a field. Let g be a Lie K-algebra.
\end{enumerate}

\begin{enumerate}
\def\labelenumi{(\alph{enumi})}
\item
  Let M and N be g-modules, where N is finite-dimensional over K. Let \((f_{j})_{1 \leqslant j \leqslant n}\) be a basis of N over K. If there are elements \(e_{1}, \ldots, e_{n}\) of M, not all zero such that \(\sum_{i=1}^{n} e_i \otimes f_i\) is invariant in \(\mathbf{M} \otimes \mathbf{N}\) , then the subspace \(\mathbf{M}_1\) of \(\mathbf{M}\) generated by the \(e_i\) is stable under the \(x_{\mathbf{M}}\) ; if further \(\mathbf{N}\) is simple the representation of \(\mathfrak{g}\) on \(\mathbf{M}_1\) is dual to the representation of \(\mathfrak{g}\) on \(\mathbf{N}\) .
\item
  Let \(M_{1}\) , \(M_{2}\) , N, \(N^{*}\) be g-modules of finite dimension over K. Suppose that \(M_{1}\) and \(M_{2}\) are simple and that the representations of g on N and \(N^{*}\) are dual. For the g-module \(M_{1}\) to be isomorphic to a sub-g-module of \(M_{2} \otimes N\) , it is necessary and sufficient that \(M_{2}\) be isomorphic to a sub-g-module of \(M_{1} \otimes N^{*}\) . (Use Proposition 4; consider the representation of g on \(M_{1}^{*} \otimes M_{2} \otimes N\) and apply (a) to the dual representation of the latter.)
\end{enumerate}

\begin{enumerate}
\def\labelenumi{\arabic{enumi}.}
\setcounter{enumi}{4}
\tightlist
\item
  Suppose that K is a field of characteristic 0. Let V be a finite-dimensional vector K-space, \(\mathfrak{g} = \mathfrak{sl}(V)\) , U the enveloping algebra of \(\mathfrak{g}\) , \(U_n\) the set of elements of filtration \(\leqslant n\) in U, \(U^n\) the set of images in U of homogeneous symmetric tensors of order \(n\) over \(\mathfrak{g}\) and \((\alpha_1, \ldots, \alpha_m)\) a basis of the vector space g. Let \(W_{s} = \bigotimes^{s} V\) . The identity representation of g defines a representation of g on \(W_{s}\) which extends to a homomorphism \(\pi_s\) of U into the algebra \(M_s = \mathcal{L}(W_s)\) .
\end{enumerate}

\begin{enumerate}
\def\labelenumi{(\alph{enumi})}
\tightlist
\item
  Let \(M_{s}^{\prime}\) be the subspace of \(M_{s}\) generated by the elements of the form \(\beta_{1} \otimes \beta_{2} \otimes \cdots \otimes \beta_{s}\) , where \(\beta_{i} \in \mathcal{L}(V)\) and \(\beta_{i} = 1\) for at least one i. Let z be an element of \(U_{s}\) and \(z_{0} = \sum_{1 \leqslant i_{1}, \ldots, i_{s} \leqslant m} a_{i_{1} \ldots i_{s}} \alpha_{i_{1}} \ldots \alpha_{i_{s}}\) its component in \(U^{s}\) , where the \(a_{i_{1} \ldots i_{s}}\) are symmetric with respect to permutation of indices. Show that
\end{enumerate}

\[
\pi_ {s} (z) \equiv s! \sum_ {1 \leqslant i _ {1}, \dots , i _ {s} \leqslant m} a _ {i _ {1} \dots i _ {s}} \alpha_ {i _ {1}} \otimes \dots \otimes \alpha_ {i _ {s}} \pmod {\mathrm{M} _ {s} ^ {\prime}}.
\]

\begin{enumerate}
\def\labelenumi{(\alph{enumi})}
\setcounter{enumi}{1}
\tightlist
\item
  Deduce from (a) that for all \(z \in \mathbf{U}\) there exists a finite-dimensional representation \(\pi\) of \(\mathbf{U}\) such that \(\pi(z) \neq 0\) . (For a sequel to this exercise, cf.~§ 7, Exercise 3.)
\end{enumerate}

\begin{enumerate}
\def\labelenumi{\arabic{enumi}.}
\setcounter{enumi}{5}
\tightlist
\item
  Suppose that \(\mathbf{K}\) is a field. Let \(\mathfrak{g}\) be a Lie K-algebra, \(x \mapsto x_{\mathbf{M}}\) a finite-dimensional representation of \(\mathfrak{g}\) , \((e_1, \ldots, e_n)\) a basis of \(\mathbf{M}\) and \((e_1^*, \ldots, e_n^*)\) the dual basis.
\end{enumerate}

\begin{enumerate}
\def\labelenumi{(\alph{enumi})}
\tightlist
\item
  If \(f\) is a linear form on \(\bigotimes^{r} \mathbf{M}\) , then:
\end{enumerate}

\[
f = \sum_ {1 \leqslant i _ {1} \dots , i _ {r} \leqslant n} f (e _ {i _ {1}} \otimes \dots \otimes e _ {i _ {r}}) (e _ {i _ {1}} ^ {*} \otimes \dots \otimes e _ {i _ {r}} ^ {*}).
\]

\begin{enumerate}
\def\labelenumi{(\alph{enumi})}
\setcounter{enumi}{1}
\tightlist
\item
  Let \(\beta\) be a non-degenerate bilinear form on \(\mathbf{M}\) which is invariant under \(\mathfrak{g}\) . Let \((e_1', \ldots, e_n')\) be the basis of \(\mathbf{M}\) such that \(\beta(e_i, e_j') = \delta_{ij}\) . Let \(f\) be an invariant linear form on \(\bigotimes^r \mathbf{M}\) . Deduce from (a) that
\end{enumerate}

\[
\sum_ {1 \leqslant i _ {1}, \dots , i _ {r} \leqslant n} f (e _ {i _ {1}} \otimes \dots \otimes e _ {i _ {r}}) (e _ {i _ {1}} ^ {\prime} \otimes \dots \otimes e _ {i _ {r}} ^ {\prime})
\]

is an invariant element of \(\otimes\) M independent of the choice of basis \((e_i)\) .

\begin{enumerate}
\def\labelenumi{(\alph{enumi})}
\setcounter{enumi}{2}
\item
  Let U be the enveloping algebra of g and U* the dual space of U. The adjoint representation of g can be extended to a representation \(x \mapsto x_{\mathrm{U}}\) of g on U defined by \(x_{\mathrm{U}}u = xu - ux\) for all \(u \in U\) . Hence U* has a g-module structure. For an element f of U* to be invariant, it is necessary and sufficient that \(f(uv) = f(vu)\) for all u, v in U.
\item
  Let \(\mathfrak{h}\) be a finite-dimensional ideal of \(\mathfrak{g}\) and \(\gamma\) an invariant bilinear form on \(\mathfrak{g}\) whose restriction to \(\mathfrak{h}\) is non-degenerate. Let \((y_i)_{1\leq i\leq n}, (y_j')_{1\leq j\leq n}\) be two bases of \(\mathfrak{h}\) such that \(\gamma(y_i,y_j') = \delta_{ij}\) . Let \(r\) be an integer \(\geqslant 1\) . Let \(f\) be a linear form on U such that \(f(w) = f(vu)\) for all \(u,v\) in U. Deduce from (b) and (c) that
\end{enumerate}

\[
\sum_ {1 \leqslant t _ {1}, \dots t _ {r} \leqslant n} f (y _ {t _ {1}} y _ {t _ {2}} \cdot \cdot \cdot y _ {t _ {r}}) y _ {t _ {1}} ^ {\prime} y _ {t _ {2}} ^ {\prime} \cdot \cdot \cdot y _ {t _ {r}} ^ {\prime}
\]

is an element of the centre of U independent of the choice of basis \((y_{i})\) . In particular recover the Casimir elements.

\begin{enumerate}
\def\labelenumi{(\alph{enumi})}
\setcounter{enumi}{4}
\tightlist
\item
  Let \(\theta\) be a permutation of \(\{1, \ldots, r\}\) . Arguing similarly, show that
\end{enumerate}

\[
\sum_ {1 \leqslant i _ {1}, \dots i _ {r} \leqslant n} f (y _ {i _ {\theta (1)}} y _ {i _ {\theta (2)}} \dots y _ {i _ {\theta (r)}}) y _ {i _ {1}} ^ {\prime} y _ {i _ {2}} ^ {\prime} \dots y _ {i _ {r}} ^ {\prime}
\]

is an element of the centre of U independent of the choice of basis \((y_{i})\) .

\begin{enumerate}
\def\labelenumi{\arabic{enumi}.}
\setcounter{enumi}{6}
\item
  Let g be a complex Lie algebra, \(\beta\) its Killing form and \(g_{0}\) the Lie algebra derived from g by restricting the base field to R. Show that the Killing form of \(g_{0}\) is twice the real part of \(\beta\) .
\item
  Let \(\mathfrak{g}\) be a real Lie algebra. The multilinear form
\end{enumerate}

\[
(x _ {1}, \dots , x _ {n}) \mapsto \operatorname{Tr} (\operatorname{ad} x _ {1} \circ \operatorname{ad} x _ {2} \circ \dots \circ \operatorname{ad} x _ {n})
\]

on \(\mathfrak{g}^n\) is invariant.

\begin{enumerate}
\def\labelenumi{\arabic{enumi}.}
\setcounter{enumi}{8}
\item
  Let g be a Lie K-algebra, \(\rho\) and \(\rho'\) semi-simple representations of g on K-modules V, \(V'\) and \(\phi\) a homomorphism of the g-module V onto the g-module \(V'\) . Show that the image under \(\phi\) of the set of invariants of V is the set of invariants of \(V'\) (use Proposition 6).
\item
  Suppose that K is a field. Let g be a Lie K-algebra and \(M_{1}\) , \(M_{2}\) two non-isomorphic simple g-modules of finite dimension over K. If \(K_{1}\) is a separable extension of K, show that there exists no simple \(g_{(K_{1})}\) -module isomorphic both to a sub- \(g_{(K_{1})}\) -module of \(M_{1(K_{1})}\) and to a sub- \(g_{(K_{1})}\) -module of \(M_{2(K_{1})}\) . (Use Proposition 4 and note that the existence of an invariant element \(\neq0\) in \(M_{1(K_{1})}^{*}\otimes_{K}M_{2(K_{1})}\) , implies the existence of an invariant element \(\neq0\) in \(M_{1}^{*}\otimes_{K}M_{2}\) (Algebra, Chapter II, § 5, no. 3, Theorem 1.)
\item
  Let \(\mathfrak{w}\) be the Lie \(p\) -algebra defined in Exercise 21 of § 1. Show that the Killing form of \(\mathfrak{w}\) is zero.
\end{enumerate}

¶ 12. Suppose that K is a field. Let g be a Lie K-algebra and M a g-module. Let \(C^{p}(\mathfrak{g}, \mathrm{M})\) denote the space of alternating linear mappings of \(g^{p}\) into M. We write \(C^{0}(\mathfrak{g}, \mathrm{M}) = \mathrm{M}\) and, for p \textless{} 0, \(C^{p}(\mathfrak{g}, \mathrm{M}) = \{0\}\) . Let \(C^{*}(\mathfrak{g}, \mathrm{M})\) be the direct sum of the \(C^{p}(\mathfrak{g}, \mathrm{M})\) . The elements of \(C^{*}(\mathfrak{g}, \mathrm{M})\) are called the cochains of g with values in M; those of \(C^{p}(\mathfrak{g}, \mathrm{M})\) are said to be of degree p.~For all \(y \in g\) , we denote by \(i(y)\) the endomorphism of \(C^{*}(\mathfrak{g}, \mathrm{M})\) which maps each subspace \(C^{p}(\mathfrak{g}, \mathrm{M})\) into \(C^{p-1}(\mathfrak{g}, \mathrm{M})\) and which, for p \textgreater{} 0, is given by the formula:

\[
(i (y) f) \left(x _ {1}, \dots , x _ {p - 1}\right) = f \left(y, x _ {1}, \dots , x _ {p - 1}\right).\tag{1}
\]

We have \(i(y)^{2}=0\) .

\begin{enumerate}
\def\labelenumi{(\alph{enumi})}
\tightlist
\item
  The adjoint representation of g and the representation of g on M define a representation of g on the space of multilinear mappings of \(g^{p}\) into M. Show that \(C^{p}(g, M)\) is stable under this representation. Let \(\theta\) be the representation of g on \(C^{*}(g, M)\) thus defined. Show that:
\end{enumerate}

\[
\theta (x) i (y) - i (y) \theta (x) = i ([ x, y ])
\]

for all x, y in g.

\begin{enumerate}
\def\labelenumi{(\alph{enumi})}
\setcounter{enumi}{1}
\tightlist
\item
  Show that there exists one and only one endomorphism \(d\) of \(\mathrm{C}^*(\mathfrak{g},\mathrm{M})\) , mapping \(\mathrm{C}^p (\mathfrak{g},\mathrm{M})\) into \(\mathrm{C}^{p + 1}(\mathfrak{g},\mathrm{M})\) , such that
\end{enumerate}

\[
d i (y) + i (y) d = \theta (y)\tag{3}
\]

for all \(y \in \mathfrak{g}\) . (Argue by induction on the degree of the cochains.) Show that, for \(f \in \mathrm{C}^p(\mathfrak{g}, \mathbf{M})\) :

\[
\begin{array}{r l} d f (x _ {1}, x _ {2}, \dots , x _ {p + 1}) & = \sum_ {i <   j} (- 1) ^ {i + j} f ([ x _ {i}, x _ {j} ], x _ {1}, \dots , \hat {x} _ {i}, \dots , \hat {x} _ {j}, \dots , x _ {p + 1}) \\ & \quad + \sum_ {i} (- 1) ^ {i + 1} (x _ {i}) _ {\mathrm{M}} f (x _ {1}, \dots , \hat {x} _ {i}, \dots , x _ {p + 1}) \end{array}
\]

(where the symbol \^{} over a letter means that it is omitted).

\begin{enumerate}
\def\labelenumi{(\alph{enumi})}
\setcounter{enumi}{2}
\tightlist
\item
  Show that, for all \(y \in g\) ,
\end{enumerate}

\[
d \theta (y) = \theta (y) d.\tag{4}
\]

(Show first, using (2) and (3), that \(d\theta(y) - \theta(y)d\) commutes with every \(i(x)\) . Then argue by induction on the degree of the cochains.)

\begin{enumerate}
\def\labelenumi{(\alph{enumi})}
\setcounter{enumi}{3}
\item
  Show that \(d^2 = 0\) . (Show first, using (3) and (4), that \(d^2\) commutes with every \(i(x)\) . The argue by induction on the degree of the cochains.)
\item
  The restriction of \(d\) to \(\mathrm{C}^p (\mathfrak{g},\mathrm{M})\) has a kernel \(\mathrm{Z}^p (\mathfrak{g},\mathrm{M})\) whose elements are called cocycles of degree \(p\) with values in M. The restriction of \(d\) to \(\mathrm{C}^{p - 1}(\mathfrak{g},\mathrm{M})\) has image \(\mathrm{B}^p (\mathfrak{g},\mathrm{M})\) , whose elements are called coboundaries of degree \(p\) with values in M. Then \(\mathrm{B}^p (\mathfrak{g},\mathrm{M})\subset \mathrm{Z}^p (\mathfrak{g},\mathrm{M})\) . The quotient space
\end{enumerate}

\[
\mathrm{Z} ^ {p} (\mathfrak {g}, \mathrm{M}) / \mathrm{B} ^ {p} (\mathfrak {g}, \mathrm{M}) = \mathrm{H} ^ {p} (\mathfrak {g}, \mathrm{M})
\]

 is called the cohomology space of g of degree p with values in M. The direct sum of the \(\mathrm{H}^{p}(\mathfrak{g},\mathrm{M})\) is denoted by \(\mathrm{H}^{*}(\mathfrak{g},\mathrm{M})\) . Show that \(\mathrm{H}^{0}(\mathfrak{g},\mathrm{M})\) is identified with the set of invariants of M. Let \(\phi\) be a homomorphism of the g-module M into the g-module N. For all \(f \in \mathrm{C}^{p}(\mathfrak{g},\mathrm{M})\) , \(\phi \circ f \in \mathrm{C}^{p}(\mathfrak{g},\mathrm{N})\) , so that \(\phi\) can be extended to a K-linear mapping \(\phi'\) : \(\mathrm{C}^{*}(\mathfrak{g},\mathrm{M}) \to \mathrm{C}^{*}(\mathfrak{g},\mathrm{N})\) . Show that \(\phi' \circ d = d \circ \phi'\) . Deduce that \(\phi'\) defines a homomorphism

\[
\tilde {\phi} \colon \mathrm{H} ^ {*} (\mathfrak {g}, \mathrm{M}) \to \mathrm{H} ^ {*} (\mathfrak {g}, \mathrm{N})
\]

which is said to be associated with \(\phi\) .

\begin{enumerate}
\def\labelenumi{(\alph{enumi})}
\setcounter{enumi}{5}
\tightlist
\item
  Let L be a sub-g-module of M and N the quotient g-module M/L. The canonical homomorphisms \(L \stackrel{\iota}{\rightarrow} M \stackrel{p}{\rightarrow} N\) define homomorphisms
\end{enumerate}

\[
\begin{array}{l} \mathrm{C} ^ {*} (\mathfrak {g}, \mathrm{L}) \xrightarrow {\iota^ {\prime}} \mathrm{C} ^ {*} (\mathfrak {g}, \mathrm{M}) \xrightarrow {p ^ {\prime}} \mathrm{C} ^ {*} (\mathfrak {g}, \mathrm{N}), \\ \mathrm{H} ^ {n} (\mathfrak {g}, \mathrm{L}) \xrightarrow {\iota^ {n}} \mathrm{H} ^ {n} (\mathfrak {g}, \mathrm{M}) \xrightarrow {p ^ {n}} \mathrm{H} ^ {n} (\mathfrak {g}, \mathrm{N}). \end{array}
\]

Show that \(\iota'\) is injective and has image the kernel of \(p'\) and that \(p'\) is surjective. For all \(c \in \mathrm{H}^n(\mathfrak{g},\mathrm{N})\) , let \(z \in \mathrm{Z}^n(\mathfrak{g},\mathrm{N})\) be a representative of \(c\) and \(a \in \mathrm{C}^n(\mathfrak{g},\mathrm{M})\) be such that \(p'(a) = z\) ; show that \(da \in \mathrm{Z}^{n+1}(\mathfrak{g},\mathrm{L})\) and that its class in \(\mathrm{H}^{n+1}(\mathfrak{g},\mathrm{L})\) depends only on \(c\) ; denote this class by \(\delta^n c\) . Show that the sequence of homomorphisms:

\[
\begin{array}{r l} 0 \longrightarrow & \mathrm{H} ^ {0} (\mathfrak {g}, \mathrm{L}) \stackrel {{\iota ^ {0}}} {{\longrightarrow}} \mathrm{H} ^ {0} (\mathfrak {g}, \mathrm{M}) \stackrel {{p ^ {0}}} {{\longrightarrow}} \mathrm{H} ^ {0} (\mathfrak {g}, \mathrm{N}) \stackrel {{\delta^ {0}}} {{\longrightarrow}} \\ & \longrightarrow \mathrm{H} ^ {1} (\mathfrak {g}, \mathrm{L}) \stackrel {{\iota ^ {1}}} {{\longrightarrow}} \mathrm{H} ^ {1} (\mathfrak {g}, \mathrm{M}) \stackrel {{p ^ {1}}} {{\longrightarrow}} \mathrm{H} ^ {1} (\mathfrak {g}, \mathrm{N}) \stackrel {{\delta^ {1}}} {{\longrightarrow}} \dots \end{array}
\]

is an exact sequence.

\begin{enumerate}
\def\labelenumi{(\alph{enumi})}
\setcounter{enumi}{6}
\tightlist
\item
  With the notation of (f), the exact sequence:
\end{enumerate}

\[
0 \to \mathscr {L} (\mathrm{N}, \mathrm{L}) \to \mathscr {L} (\mathrm{N}, \mathrm{M}) \to \mathscr {L} (\mathrm{N}, \mathrm{N}) \to 0
\]

defines an exact sequence

\[
\mathrm{H} ^ {0} (\mathfrak {g}, \mathscr {L} (\mathrm{N}, \mathrm{M})) \rightarrow \mathrm{H} ^ {0} (\mathfrak {g}, \mathscr {L} (\mathrm{N}, \mathrm{N})) \rightarrow \mathrm{H} ^ {1} (\mathfrak {g}, \mathscr {L} (\mathrm{N}, \mathrm{L})).
\]

 The identity mapping of N is an invariant u of \(\mathcal{L}(\mathrm{N},\mathrm{N})\) and hence an element of \(\mathrm{H}^{0}(\mathfrak{g},\mathcal{L}(\mathrm{N},\mathrm{N}))\) ; let c be its image in \(\mathrm{H}^{1}(\mathfrak{g},\mathcal{L}(\mathrm{N},\mathrm{L}))\) . Then for the existence in M of a sub-g-module supplementary to L it is necessary and sufficient that c=0. (This condition means that u is the image of an element of \(\mathrm{H}^{0}(\mathfrak{g},\mathcal{L}(\mathrm{N},\mathrm{M}))\) , that is a homomorphism v of the g-module N into the g-module M; then \(v(\mathrm{N})\) is the desired supplement.)

\begin{enumerate}
\def\labelenumi{(\alph{enumi})}
\setcounter{enumi}{7}
\item
  Show that \(Z^{1}(\mathfrak{g},\mathfrak{g})\) (where g is considered as a g-module by means of the adjoint representation) is identified with the vector space of derivations of g and that \(B^{1}(\mathfrak{g},\mathfrak{g})\) is identified with the vector space of inner derivations of g.
\item
  Let \(\mathfrak{a}\) and \(\mathfrak{b}\) be two Lie K-algebras with \(\mathfrak{b}\) commutative and \(\mathfrak{b} \xrightarrow{\lambda} \mathfrak{g} \xrightarrow{\mu} \mathfrak{a}\) an extension of \(\mathfrak{a}\) by \(\mathfrak{b}\) . For all \(x \in \mathfrak{g}\) , the restriction of \(\mathrm{ad}_{\mathfrak{a}} x\) to \(\mathfrak{b}\) depends only on the class of \(x\) modulo \(\mathfrak{b}\) , whence there is an \(\mathfrak{a}\) -module structure on \(\mathfrak{b}\) . Let \(\nu\) be a K-linear mapping of \(\mathfrak{a}\) into \(\mathfrak{g}\) such that \(\mu \circ \nu\) is the identity mapping of \(\mathfrak{a}\) . For \(x, y\) in \(\mathfrak{a}\) , we write \(f(x, y) = [\nu x, \nu y] - \nu([x, y])\) . Show that \(f \in Z^2(\mathfrak{a}, \mathfrak{b})\) and that the class \(c\) of \(f\) in \(H^2(\mathfrak{a}, \mathfrak{b})\) does not depend on the choice of \(\nu\) . For two extensions of \(\mathfrak{a}\) by \(\mathfrak{b}\) , which define the same \(\mathfrak{a}\) -module structure on \(\mathfrak{b}\) , to be equivalent, it is necessary and sufficient that the corresponding class \(c \in H^2(\mathfrak{a}, \mathfrak{b})\) be the same. For the extension to be inessential, it is necessary and sufficient that \(c = 0\) . If B is an \(\mathfrak{a}\) -module and \(c \in H^2(\mathfrak{a}, \mathfrak{B})\) , there exists an extension of \(\mathfrak{a}\) by B (considered as a commutative Lie algebra) defining the given \(\mathfrak{a}\) -module structure on B and the given element \(c\) of \(H^2(\mathfrak{a}, \mathfrak{B})\) .
\item
  Let g be a finite-dimensional Lie algebra over K, U its enveloping algebra, h an ideal of g, β an invariant bilinear form on g whose restriction to h is non-degenerate, \((y_{i})_{1 \leqslant i \leqslant n}\) and \((y_{i}')_{1 \leqslant i \leqslant n}\) two bases of h such that
\end{enumerate}

 \(\beta(y_i, y_i') = \delta_{ij}, t = \sum_{i=1}^{n} y_i y_i' \in \mathbf{U}, \mathbf{M}\) a \(g\) -module and \(\rho\) the endomorphism of \(\mathrm{C}^*(\mathfrak{g}, \mathbf{M})\) which maps each \(\mathrm{C}^p(\mathfrak{g}, \mathbf{M})\) into \(\mathrm{C}^{p-1}(\mathfrak{g}, \mathbf{M})\) and which, for \(p > 0\) , is defined by:

\[
(\wp f) (x _ {1}, \dots , x _ {p - 1}) = \sum_ {k = 1} ^ {n} (y _ {k}) _ {\mathrm{M}} f (y _ {k} ^ {\prime}, x _ {1}, \dots , x _ {p - 1}).
\]

Finally let \(\Gamma\) be the endomorphism of \(C^*(g, M)\) extending \(t_M\) which maps each cochain \(f\) of degree \(>0\) onto the cochain \(t_M \circ f\) . Show that \(\rho d + d\rho = \Gamma\) . (If for \(x \in g\) we write:

\[
[ x, y _ {k} ] = \sum_ {l = 1} ^ {n} a _ {k l} y _ {l}, \quad [ x, y _ {k} ^ {\prime} ] = \sum_ {l = 1} ^ {n} a _ {k l} ^ {\prime} y _ {l} ^ {\prime},
\]

·show first that \(a_{kl} = -a_{lk}^{\prime}\) .) Deduce that \(\Gamma d = d\Gamma\) . If M is simple and finite-dimensional over K, \(\beta\) is the associated bilinear form and dim \(\mathfrak{h}\) is not divisible by the characteristic of K, show that \(\mathrm{H}^{p}(\mathfrak{g},\mathrm{M})=\{0\}\) for all p.~(Using Proposition 12, show that \(\Gamma\) is an automorphism of \(\mathrm{C}^{*}(\mathfrak{g},\mathrm{M})\) and hence induces automorphisms of \(\mathrm{Z}^{p}(\mathfrak{g},\mathrm{M})\) and \(\mathrm{B}^{p}(\mathfrak{g},\mathrm{M})\) . For

\[
f \in \mathrm{Z} ^ {p} (\mathfrak {g}, \mathrm{M}), \quad \Gamma f = (d \rho + \rho d) f = d \rho f \in \mathrm{B} ^ {p} (\mathfrak {g}, \mathrm{M}),
\]

whence \(Z^{p}(\mathfrak{g},\mathbf{M}) = \mathrm{B}^{p}(\mathfrak{g},\mathbf{M}).\)

